\section{Kinematics}\label{sec_kinematics}
Following the standard David-Hartenberg convention for assignment of local coordinate systems, it is possible to derive the David-Hartenberg parameters. The local coordinate systems are shown on figure \ref{fig_cobra_dh}.
\begin{figure}[htbp]
\centering
\begin{tikzpicture}
\node[anchor=south west,inner sep=0] (image) at (0,0) 
{\includegraphics[width=0.6\textwidth]{kinematics_fig}};
    \begin{scope}[x={(image.south east)},y={(image.north west)}]
        %Rotation axis
        \draw[blue,dashed, ultra thick] (0.73,-0.1) -- (0.73,0.7);
        \draw[blue,dashed,ultra thick] (0.37,0.3) -- (0.37,0.7);
        \draw[blue,dashed,ultra thick] (0.055,0.12) -- (0.055,0.7);
        
        % Local coordinates
        
        % q0
        \draw[red,ultra thick, ->] (0.73,0) --node[right]{$\boldsymbol{z_0}$} (0.73,0.1);
        \draw[red,ultra thick, ->] (0.73,0) --node[below]{$\boldsymbol{x_0}$} (0.63,0);
        
        % q1
        \draw[red,ultra thick, ->] (0.37,0.5) --node[right]{$\boldsymbol{z_1}$} (0.37,0.6);
        \draw[red,ultra thick, ->] (0.37,0.5) --node[below]{$\boldsymbol{x_1}$} (0.27,0.5);
        
        %q2
        \draw[red,ultra thick, ->] (0.055,0.5) --node[right]{$\boldsymbol{z_2}$} +(0,0.1);
        \draw[red,ultra thick, ->] (0.055,0.5) -- node[below]{$\boldsymbol{x_2}$}+(-0.1,0); 
        
        %q3
        \draw[red,ultra thick, ->] (0.055,0.19) --node[right,thick]{$\boldsymbol{z_3}$} (0.055,0.29); 
        \draw[red,ultra thick, ->] (0.055,0.19) --node[below]{$\boldsymbol{x_3}$} +(-0.1,0);
        
        %q4
        \draw[red,ultra thick, ->] (0.055,0.13) --node[right,thick]{$\boldsymbol{z_4}$} +(0,-0.1);
        \draw[red,ultra thick, ->] (0.055,0.13) --node[below]{$\boldsymbol{x_4}$} +(-0.1,0);
        
        % measurements
        \draw [green, ultra thick,|-|] (0.73,0.35) --node[below]{$\boldsymbol{a_1}$} (0.37,0.35);
        \draw [green, ultra thick,-|] (0.37,0.35) --node[below]{$\boldsymbol{a_2}$} (0.055,0.35);
        \draw [green, ultra thick,|-|] (0.83,0) --node[right]{$\boldsymbol{d_1}$} (0.83,0.5);
        \draw [green, ultra thick,|-|] (-0.15,0.5) --node[left]{$\boldsymbol{q_3}$} (-0.15,0.19);
         \draw [green, ultra thick,-|] (-0.15,0.19) -- node[left]{$\boldsymbol{d_4}$} (-0.15,0.13);
    \end{scope}
\end{tikzpicture}

\caption{The Cobra s600 with local coordinate systems.}
\label{fig_cobra_dh}
\end{figure}
From the figure it is possible to setup the DH-table, which is shown in table \ref{tab_dh}
\begin{table}[htbp!]
\centering
  \begin{tabular}{c|c|c|c|c}
  $i$ & $\alpha_i$ & $a_i$ & $d_i$ & $\theta_i$ \\\hline
  1 & 0 & $a_1$ & $d_1$ & $q_1$ \\\hline
  2 & 0 & $a_2$ & 0 & $q_2$ \\\hline
  3 & 0 & 0 & $q_3$ & 0 \\\hline
  4 & $\pi$ & 0 & $d_4$ & $q_4$ \\
    
  \end{tabular}
  \caption{The David-Hartenberg parameters for the Cobra s600}
  \label{tab_dh}
\end{table}\newline
The general DH-matrix is as follows when using compact notation:
\begin{equation}
  ^{i-1}A_i (q_i)=\mat{
  c\theta_i & -c\alpha_i\g s\theta_i & s\alpha_i\g s\theta_i & a_i\g c\theta_i \\
  s\theta_i & c\alpha_i\g c\theta_i & -s\alpha_i \g c\theta_i & a_i\g s\theta_i \\
  0 & s\alpha_i & c\alpha_i & d_i\\
  0 & 0 & 0 & 1}
\end{equation}
To derive the transformation matrix from the tool tip to the Cartesian space, it is necessary to multiply the four transformation matrices. This is done in appendix \ref{app_dh}. The final transformation matrix is:
\begin{equation}
  \T{0}{4}=
  \mat{c_{124} & s_{124}    & 0   & a_1\g c_1+a_2\g c_{12}\\
       s_{124} & -c_{124}   & 0   & a_1\g s_1+a_2\g s_{12}\\
       0       &   0        & -1  & d_1+q_3+d_4\\
       0       &   0        & 0   & 1}\label{eq_t4}
\end{equation}

